\section{Results that argue against the method}
\label{sec:against}

Five rows are reported unchanged, and three of them are outright negative ---
worse than a baseline that needs no training at all.

The London office returns skill $-0.463$ and public services in Washington DC
returns $+0.052$. Both sit at $1$--$3\%$ HVAC share, where a weather-aware model
has little to add over ``last week, same time'' --- a coherent and testable
explanation rather than an excuse. Delhi returns $+0.097$ with the study's worst
metered calibration, and its explanation is different and given in
Section~\ref{sec:india-fix}: a system-level series carrying calendar structure
the feature set does not encode. Guangdong and Heilongjiang return $-0.314$ and
$-0.294$, the two lowest figures in the study, and Section~\ref{sec:china}
explains why those two should not be used against the method any more than they
should have been used for it.

Two defects in our own pipeline were exposed by running it outside its home
site. A hard-coded United States federal holiday calendar would have inverted
the holiday flag on two of the largest annual load anomalies for every European
building; holidays are now resolved per country, and India is the case that
forces it, its calendar being substantially lunisolar and therefore not
expressible as a fixed list. And a data-preparation stage overwrote its own
manifest, silently discarding every site but the last, which single-site work
had never exposed.

A third is worth recording because it is the kind that survives review. The
column commonly written as ``the Boole bound'' in applied chance-constraint work
is usually $1-(1-\alpha)^H$, which is the independence calculation and is not a
bound. Our own tables carried that label until this draft. The corrected reading
in Section~\ref{sec:horizon} is stronger than the one it replaces, because the
valid bound turns out to be vacuous over the horizon that matters.
